% SOURCE_AUTHORITY: sga5_fr_workpass.tex, lines 13399--13439 (LF-normalized unit).
% SOURCE_SHA256: 791F4EFFC5E02832D5D77ED03518C8156D6F07E4C8238B03545DB93D883FBB28
\section*{4. Generalidades sobre os termos locais}

Suponhamos que $X'$ seja liso sobre $k$ e, além disso, que o número de pontos fixos de $f$ seja finito. Daí resulta, em particular, que o número de pontos fixos de $g^{-1}f^{\prime}$ é finito para cada $g \in G$.

Para cada ponto fixo $x \in X$ de $f$, introduz-se a função $\varphi$-central
\[
S_x(f',f,G,\varphi) : G \longrightarrow \bQ \quad ,
\]
com valores racionais, definida pela fórmula:

\begin{equation}\tag{4.1}
S_x(f',f,G,\varphi)(g) = \frac{1}{c^\varphi(g)} \sum_{x' \in {X'_x}^{g^{-1}f'}} (\nu_{x'}(g^{-1}f'_R) - 1) \quad .
\end{equation}

As seguintes notações foram utilizadas:

\begin{equation}\tag{4.2}
c^\varphi(g) = \card G_g^\varphi \quad , \quad G_g^\varphi = \{ x \in G / \varphi(x)gx^{-1} = g \} \quad ,
\end{equation}

\begin{equation}\tag{4.3}
{X'_x}^{g^{-1}f'} = \{ x' \in X' / p(x') = x \; , \; (g^{-1}f^{\prime})(x') = x' \} \quad .
\end{equation}

\begin{equation}\tag{4.4}
\nu_{x'}(g^{-1}f') = \text{a multiplicidade do ponto fixo } x' \text{ de } g^{-1}f^{\prime} \quad .
\end{equation}

Por composição com as inclusões canônicas
\[
\bQ \hookrightarrow \bQ_\ell
\]
A função $S_x(f',f,G,\varphi)$ dá origem a funções $\varphi$-centrais $S_{x,\ell}(f',f,G,\varphi)$ sobre $G$ com valores no corpo $\bQ_\ell$ dos números $\ell$-ádicos.

\medskip
\noindent
\textbf{Conjectura 4.5.} \emph{Para todo $\ell \neq \car(k)$, os valores da função $S_{x,\ell}(f',f,G,\varphi)$ são inteiros $\ell$-ádicos.}

\medskip

A seguir, demonstraremos que a conjectura 4.5 é verdadeira quando a dimensão de $X$ é igual a 1.
